\section{Results}\label{sec:results}

\todoblock{This section is intentionally a scaffold. Insert results
\emph{verbatim} from the aggregation CSVs once
\texttt{configs/article\_3vrf\_xgb\_magic.yaml} (and its multi-dataset
variant to be created) finishes running on the v1 panel of
Section~\ref{sec:experimental-design}. Do not invent numbers.}

\subsection{Dataset-level performance}
\label{sec:res-performance}
\begin{table}[t]
  \centering
  \caption{Predictive performance (mean $\pm$ standard deviation over
  $R{=}10$ replicas) per (dataset, algorithm, method).
  \textbf{TODO: insert after experiment runs.}}
  \label{tab:performance}
  \begin{tabular}{lllrrrrr}
    \toprule
    Dataset & Algorithm & Method & ROC-AUC & PR-AUC & F1 & Bal.\ Acc. & MCC \\
    \midrule
    \multicolumn{8}{c}{\textit{TODO: rows from
       experiments/<dataset>/<design>/aggregated\_table.csv}} \\
    \bottomrule
  \end{tabular}
\end{table}

\subsection{Computational cost}
\label{sec:res-cost}
\begin{table}[t]
  \centering
  \caption{Computational cost: mean time per fold, total optimization
  time, and total wall-clock time. \textbf{TODO}.}
  \label{tab:cost}
  \begin{tabular}{lllrrr}
    \toprule
    Dataset & Algorithm & Method & Time/fold (s) & Opt.\ time (s) & Total time (s) \\
    \midrule
    \multicolumn{6}{c}{\textit{TODO}} \\
    \bottomrule
  \end{tabular}
\end{table}

\subsection{Cross-dataset comparison}
\label{sec:res-friedman}
\todoblock{Insert Friedman test and Nemenyi critical-difference plot
after running the panel. Discuss any rank inversions between the
proposed method and TPE/BO. Cite \citet{demsar2006statistical}.}

\subsection{NBI vs.\ legacy weighted-sum ablation}
\label{sec:res-ablation}
\todoblock{Side-by-side comparison: legacy
\texttt{run\_nbi\_weighted\_sum} vs.\ true $N$-objective NBI. The
expected outcome is that NBI produces more uniformly spaced points on
the Pareto front, but the practical performance difference depends on
front curvature; a synthetic confirmation already passes in
\texttt{tests/methodology/test\_nbi\_core.py}.}

\subsection{$N$-objective NBI residual diagnostics}
\label{sec:res-nbi-residuals}
\begin{table}[t]
  \centering
  \caption{NBI sub-problem residual norms across the $51$ $\beta$ /
  $66$ $\beta$ grid for the $q{=}2$ / $q{=}3$ runs, per dataset.
  \textbf{TODO}.}
  \label{tab:nbi-residuals}
  \begin{tabular}{lrrr}
    \toprule
    Dataset & $\max \|\text{residual}\|$ & median & $p_{95}$ \\
    \midrule
    \multicolumn{4}{c}{\textit{TODO}} \\
    \bottomrule
  \end{tabular}
\end{table}

\subsection{Factor model and VRF interpretation}
\label{sec:res-vrf}
\todoblock{Insert per-dataset rotated loading tables (one row per raw
metric, one column per VRF), with eigenvalues and cumulative
variance. Highlight the construct each VRF maps to.}

\subsection{MBPA post-optimization trigger summary}
\label{sec:res-mbpa}
\begin{table}[t]
  \centering
  \caption{Per-dataset MBPA trigger summary. ``Triggered'' = at least
  one of the six diagnostics fired; ``Refined'' = a refined
  $\bm w^{*}$ was recovered; ``Improved'' = the refined candidate
  improved the selection metric over the primary $\bm w^{*}$.
  \textbf{TODO}.}
  \label{tab:mbpa}
  \begin{tabular}{lrrrr}
    \toprule
    Dataset & Replicas & Triggered & Refined & Improved \\
    \midrule
    \multicolumn{5}{c}{\textit{TODO}} \\
    \bottomrule
  \end{tabular}
\end{table}

\subsection{Reproducibility artifact summary}
\label{sec:res-artifacts}
\todoblock{Insert artifact count, total storage, hash of the canonical
design, and the GitHub release tag (\texttt{v0.3.0-articlev1}) once
the campaign is archived.}
